% ============================================================
\section{Day 5}
% ============================================================

% For the coupled-oscillator figure:
% \usetikzlibrary{decorations.pathmorphing}

% ============================================================
\subsection{Readings and resources}
% ============================================================

\subsubsection{Readings}

\begin{itemize}
    \item Christine Rosen, \textit{The Extinction of Experience: Being Human in a Disembodied World}.
    \item Frederick Douglass, \textit{Narrative of the Life of Frederick Douglass, an American Slave}.
\end{itemize}

\subsubsection{Popular biology books for physicists}

\begin{itemize}
    \item Ron Milo and Rob Phillips, \textit{Cell Biology by the Numbers}.
    \item Bruce Alberts et al., \textit{Essential Cell Biology}.
    \item Sean B. Carroll, \textit{The Serengeti Rules}.
    \item Mark Ptashne, \textit{A Genetic Switch}.
    \item Horace Freeland Judson, \textit{The Eighth Day of Creation}.
    \item Philip Ball, \textit{How Life Works}.
    \item Paul Nurse, \textit{What Is Life?}
    \item Ed Yong, \textit{An Immense World}.
\end{itemize}

\subsubsection{Popular physics books for biologists}

\begin{itemize}
    \item Albert Einstein and Leopold Infeld, \textit{The Evolution of Physics}.
    \item George Gamow, \textit{Biography of Physics}.
    \item I. Bernard Cohen, \textit{The Birth of a New Physics}.
\end{itemize}

\subsubsection{Mathematics resources}

\begin{itemize}
    \item Gilbert Strang, \textit{Introduction to Linear Algebra}.
    \item 3Blue1Brown, particularly the \textit{Essence of Linear Algebra} series.
\end{itemize}


% ============================================================
\subsection{Representations}
% ============================================================

The way in which a problem is represented can determine which aspects of the problem become easy or difficult to see.
A useful representation exposes the structure that matters while suppressing details that are irrelevant to the question being asked.

The same idea applies to how we describe ourselves.
Labels such as ``physicist'' or ``biologist'' are only one possible representation and may hide much of the information that actually characterizes a person.


% ------------------------------------------------------------
\subsubsection{1. Vectors}
% ------------------------------------------------------------

Many complicated objects can be represented as vectors.
Once an object has been written as a vector, the tools of linear algebra can be used to compare objects, transform representations, and identify important directions of variation.

Examples include:

\begin{itemize}
    \item an image, represented by a vector of pixel intensities;
    \item a transcriptome, represented by a vector of gene-expression levels;
    \item a biological shape, represented by coordinates or by coefficients in a suitable basis, such as spherical harmonics or finite elements;
    \item a probability distribution, represented by the probabilities associated with each possible state.
\end{itemize}


% ------------------------------------------------------------
\paragraph{Example: the Buri experiment}
% ------------------------------------------------------------

Consider again the Buri experiment.
Each vial contains 16 diploid flies, so there are

\begin{equation}
    2\times16=32
\end{equation}

copies of the relevant locus.

Let $n_R$ denote the number of red alleles in a vial.
The possible values are

\begin{equation}
    n_R=0,1,2,\ldots,32.
\end{equation}

There are therefore 33 possible allele-count states.

The allele count $n_R$ itself is a single discrete variable.
However, a probability distribution over all possible allele counts can be represented by a vector with 33 components,

\begin{equation}
    \mathbf{p} =
    \begin{pmatrix}
        p_0 \\
        p_1 \\
        \vdots \\
        p_{32}
    \end{pmatrix}.
\end{equation}

Here

\begin{equation}
    p_i=P(n_R=i).
\end{equation}

The vector therefore lives in a 33-dimensional vector space, although not every vector in that space is a valid probability distribution.
The components must satisfy

\begin{equation}
    p_i\geq0
\end{equation}

and

\begin{equation}
    \sum_{i=0}^{32}p_i=1.
\end{equation}

Thus a stochastic evolutionary process can be viewed as the time evolution of a probability vector over the possible allele-count states.


% ============================================================
\subsection{2. Fitting data}
% ============================================================

Suppose we have $N$ data points

\begin{equation}
    (x_1,y_1),(x_2,y_2),\ldots,(x_N,y_N)
\end{equation}

and wish to represent them by a straight line,

\begin{equation}
    y=ax+b.
\end{equation}

For each data point, the residual between the observed and predicted value is

\begin{equation}
    r_i=y_i-(ax_i+b).
\end{equation}

The least-squares method chooses $a$ and $b$ to minimize the sum of the squared residuals,

\begin{equation}
    E(a,b)=\sum_{i=1}^{N}\left[y_i-(ax_i+b)\right]^2.
\end{equation}

This provides a quantitative definition of the ``best-fitting'' line.


% ------------------------------------------------------------
\subsubsection{Why squares?}
% ------------------------------------------------------------

One reason squared errors arise naturally is that measurement errors are often modeled as independent Gaussian random variables.

Suppose

\begin{equation}
    y_i=ax_i+b+\epsilon_i,
\end{equation}

where

\begin{equation}
    \epsilon_i\sim\mathcal{N}(0,\sigma^2).
\end{equation}

The probability density of observing $y_i$ for given $a$ and $b$ is then

\begin{equation}
    p(y_i\mid x_i,a,b)=\frac{1}{\sqrt{2\pi\sigma^2}}\exp\left[-\frac{\left(y_i-ax_i-b\right)^2}{2\sigma^2}\right].
\end{equation}

Assuming that the errors of different measurements are independent, the likelihood of the complete dataset is

\begin{equation}
    \mathcal{L}(a,b)=\prod_{i=1}^{N}p(y_i\mid x_i,a,b).
\end{equation}

Taking the logarithm converts the product into a sum,

\begin{equation}
    \ln\mathcal{L}(a,b)=\text{constant}-\frac{1}{2\sigma^2}\sum_{i=1}^{N}\left[y_i-(ax_i+b)\right]^2.
\end{equation}

Maximizing the likelihood is therefore equivalent to minimizing

\begin{equation}
    \sum_{i=1}^{N}\left[y_i-(ax_i+b)\right]^2.
\end{equation}

Thus least squares follows naturally from a model in which the measurement noise is independent and Gaussian with constant variance.


% ------------------------------------------------------------
\subsubsection{Finding the best-fit parameters}
% ------------------------------------------------------------

At the minimum of the error function,

\begin{equation}
    \frac{\partial E}{\partial a}=0
\end{equation}

and

\begin{equation}
    \frac{\partial E}{\partial b}=0.
\end{equation}

Differentiating with respect to $a$ gives

\begin{equation}
    \frac{\partial E}{\partial a}=\sum_{i=1}^{N}-2x_i\left[y_i-(ax_i+b)\right]=0.
\end{equation}

Therefore,

\begin{equation}
    a\sum_{i=1}^{N}x_i^2+b\sum_{i=1}^{N}x_i=\sum_{i=1}^{N}x_iy_i.
\end{equation}

Differentiating with respect to $b$ gives

\begin{equation}
    \frac{\partial E}{\partial b}=\sum_{i=1}^{N}-2\left[y_i-(ax_i+b)\right]=0.
\end{equation}

Therefore,

\begin{equation}
    a\sum_{i=1}^{N}x_i+Nb=\sum_{i=1}^{N}y_i.
\end{equation}

The fitting problem has now become a system of two linear equations for the two unknown parameters $a$ and $b$.


% ------------------------------------------------------------
\subsubsection{Matrix representation}
% ------------------------------------------------------------

The two equations can be written as

\begin{equation}
    \begin{pmatrix}
        \sum_i x_i^2 & \sum_i x_i \\
        \sum_i x_i & N
    \end{pmatrix}
    \begin{pmatrix}
        a \\
        b
    \end{pmatrix}
    =
    \begin{pmatrix}
        \sum_i x_i y_i \\
        \sum_i y_i
    \end{pmatrix}.
\end{equation}

A general system of linear equations has the form

\begin{equation}
    A\mathbf{x}=\mathbf{c}.
\end{equation}

Formally, if $A$ is invertible,

\begin{equation}
    \mathbf{x}=A^{-1}\mathbf{c}.
\end{equation}

For the straight-line fit, define

\begin{equation}
    \Delta=N\sum_i x_i^2-\left(\sum_i x_i\right)^2.
\end{equation}

The slope is

\begin{equation}
    \boxed{a=\frac{N\sum_i x_i y_i-\left(\sum_i x_i\right)\left(\sum_i y_i\right)}{\Delta}.}
\end{equation}

The intercept is

\begin{equation}
    \boxed{b=\frac{\left(\sum_i x_i^2\right)\left(\sum_i y_i\right)-\left(\sum_i x_i\right)\left(\sum_i x_i y_i\right)}{\Delta}.}
\end{equation}


% ============================================================
\subsection{3. Finding the right coordinates}
% ============================================================

A representation can be formally correct and still be inconvenient.
An important strategy in physics is therefore to search for coordinates in which the important structure of the problem becomes simple.


% ------------------------------------------------------------
\subsubsection{Coupled oscillators and normal modes}
% ------------------------------------------------------------

Consider two masses coupled by springs.

\begin{figure}[h!]
\centering
\begin{tikzpicture}

    \draw[thick] (0,0) -- (0,1.2);
    \draw[thick] (8,0) -- (8,1.2);

    \draw[decorate,decoration={zigzag,segment length=5pt,amplitude=3pt}] (0,0.6) -- (2,0.6);

    \draw[thick] (2,0.2) rectangle (3,1.0);
    \node at (2.5,0.6) {$m$};

    \draw[decorate,decoration={zigzag,segment length=5pt,amplitude=3pt}] (3,0.6) -- (5,0.6);

    \draw[thick] (5,0.2) rectangle (6,1.0);
    \node at (5.5,0.6) {$m$};

    \draw[decorate,decoration={zigzag,segment length=5pt,amplitude=3pt}] (6,0.6) -- (8,0.6);

    \node below at (2.5,0.15) {$x_1$};
    \node below at (5.5,0.15) {$x_2$};

\end{tikzpicture}
\caption{Two coupled oscillators.
The individual displacements $x_1$ and $x_2$ provide one coordinate system, while the normal modes provide another.}
\end{figure}

The individual mass displacements $x_1$ and $x_2$ are natural coordinates, but their equations of motion are coupled.

For a symmetric system, more useful coordinates are the combinations

\begin{equation}
    q_{+}=\frac{x_1+x_2}{\sqrt{2}}
\end{equation}

and

\begin{equation}
    q_{-}=\frac{x_1-x_2}{\sqrt{2}}.
\end{equation}

The coordinate $q_{+}$ describes the mode in which the masses move together.
The coordinate $q_{-}$ describes the mode in which they move oppositely.

These are the normal modes of the system.
In normal-mode coordinates, the coupled problem can be decomposed into independent motions.

This illustrates the importance of finding coordinates adapted to the structure of the problem.


% ============================================================
\subsubsection{Covariance}
% ============================================================

The variance of a variable $x$ measures its fluctuations around its mean,

\begin{equation}
    \operatorname{Var}(x)=\left\langle\left(x-\langle x\rangle\right)^2\right\rangle.
\end{equation}

Expanding the square gives

\begin{equation}
    \operatorname{Var}(x)=\langle x^2\rangle-\langle x\rangle^2.
\end{equation}

For two variables $x_1$ and $x_2$, the covariance is

\begin{equation}
    \operatorname{Cov}(x_1,x_2)=\left\langle\left(x_1-\langle x_1\rangle\right)\left(x_2-\langle x_2\rangle\right)\right\rangle.
\end{equation}

The covariance matrix for two variables is

\begin{equation}
    C=
    \begin{pmatrix}
        \left\langle\left(x_1-\langle x_1\rangle\right)^2\right\rangle &
        \left\langle\left(x_1-\langle x_1\rangle\right)\left(x_2-\langle x_2\rangle\right)\right\rangle \\
        \left\langle\left(x_1-\langle x_1\rangle\right)\left(x_2-\langle x_2\rangle\right)\right\rangle &
        \left\langle\left(x_2-\langle x_2\rangle\right)^2\right\rangle
    \end{pmatrix}.
\end{equation}

The diagonal entries describe the variance of each coordinate.
The off-diagonal entries describe how fluctuations in different coordinates are correlated.


% ============================================================
\subsubsection{Principal component analysis}
% ============================================================

Principal component analysis, or PCA, provides a systematic way to find coordinates that capture the dominant variation in a dataset.

For variables collected into a vector $\mathbf{x}$, define the covariance matrix

\begin{equation}
    C_{ij}=\left\langle\left(x_i-\langle x_i\rangle\right)\left(x_j-\langle x_j\rangle\right)\right\rangle.
\end{equation}

The principal components are the eigenvectors of this matrix,

\begin{equation}
    C\mathbf{v}_{\mu}=\lambda_{\mu}\mathbf{v}_{\mu}.
\end{equation}

The eigenvectors $\mathbf{v}_{\mu}$ define new coordinate directions.
The corresponding eigenvalues $\lambda_{\mu}$ measure the variance of the data along those directions.

The first principal component is the direction along which the data vary most.
The second is the orthogonal direction containing the next-largest variance, and so on.

If only a few eigenvalues are large, a high-dimensional dataset can be represented accurately using only a few principal components.
PCA therefore provides both a change of coordinates and a method for dimensionality reduction.


% ============================================================
\subsection{The eigenworm}
% ============================================================

A biological example of finding useful coordinates is the analysis of \textit{Caenorhabditis elegans} posture and locomotion.

A related historical reading is Charles Darwin's \textit{The Formation of Vegetable Mould through the Action of Worms}.

A quantitative study of worm posture is:

\begin{quote}
G. J. Stephens, B. Johnson-Kerner, W. Bialek, and W. S. Ryu,
\textit{Dimensionality and Dynamics in the Behavior of C. elegans},
PLoS Computational Biology \textbf{4}, e1000028 (2008).
\end{quote}

The goal is to determine whether the apparently complicated set of worm postures can be represented using only a few collective coordinates.

A simplified image-analysis and PCA pipeline is:

\begin{enumerate}[start=0]
    \item Parse the image data.
    \item Segment the worm from the image background.
    \item Skeletonize the worm to determine its centerline.
    \item Fit or sample a smooth contour along the centerline.
    \item Compute the tangent angle along the body for every recorded time.
    \item Arrange the tangent-angle data as a function of body position and time.
    \item Compute the covariance matrix of the angular fluctuations.
    \item Perform principal component analysis.
    \item Decompose each observed posture into the resulting eigenmodes.
\end{enumerate}

Let $\theta(s,t)$ denote the tangent angle at position $s$ along the worm body at time $t$.
After removing the overall orientation of the worm, a posture can be represented by the set of angles along its body.

The covariance between two body positions $s$ and $s'$ is

\begin{equation}
    C(s,s')=\left\langle\left[\theta(s)-\langle\theta(s)\rangle\right]\left[\theta(s')-\langle\theta(s')\rangle\right]\right\rangle.
\end{equation}

The eigenvectors of this covariance matrix satisfy

\begin{equation}
    \sum_{s'}C(s,s')u_{\mu}(s')=\lambda_{\mu}u_{\mu}(s).
\end{equation}

These principal components are called \emph{eigenworms}.

The worm shape can then be represented approximately as

\begin{equation}
    \theta(s,t)\approx\sum_{\mu=1}^{K}a_{\mu}(t)u_{\mu}(s).
\end{equation}

Stephens et al.\ found that only four eigenworms account for more than $95\%$ of the observed variance in worm shapes.

Thus an object initially represented using roughly one hundred local angular coordinates can be represented, to good approximation, by only four mode amplitudes,

\begin{equation}
    \mathbf{a}(t)=
    \begin{pmatrix}
        a_1(t) \\
        a_2(t) \\
        a_3(t) \\
        a_4(t)
    \end{pmatrix}.
\end{equation}

The important result is not merely data compression.
The new coordinates provide a compact description of the behavior and dynamics of the animal.


% ------------------------------------------------------------
\subsubsection{Geometric phase and locomotion}
% ------------------------------------------------------------

A related example is:

\begin{quote}
J. M. Rieser et al.,
\textit{Geometric phase predicts locomotion performance in undulating living systems across scales},
PNAS \textbf{121}, e2320517121 (2024).
\end{quote}

Periodic locomotion can be represented as a closed trajectory through a low-dimensional space of body shapes.
The coordinates of this space describe how much of different shape modes are present at a given time.

As the animal cycles through a sequence of shapes, it traces a loop through shape space.
The geometry of this loop can be related to the net displacement accumulated during one locomotor cycle.

Thus an appropriate low-dimensional representation can expose a direct relationship between body-shape dynamics and locomotion.


% ============================================================
\subsection{Shannon entropy}
% ============================================================

Shannon entropy provides a measure of the uncertainty associated with a probability distribution.

For a discrete probability distribution

\begin{equation}
    \{p_1,p_2,\ldots,p_n\},
\end{equation}

the Shannon entropy is

\begin{equation}
    \boxed{S(\{p_i\})=-\sum_{i=1}^{n}p_i\ln p_i.}
\end{equation}

A distribution concentrated entirely in one state has no uncertainty and therefore zero entropy.
A broader distribution represents greater uncertainty.

When no constraints other than normalization are imposed, entropy is maximized by assigning equal probability to every available state.
This motivates the maximum-entropy principle: choose the probability distribution that satisfies the known constraints while otherwise introducing as little additional structure as possible.

In this sense, the maximum-entropy distribution is maximally non-committal about information that has not been specified.


% ------------------------------------------------------------
\subsubsection{Maximum entropy with constraints}
% ------------------------------------------------------------

Suppose that the states are labeled

\begin{equation}
    i=1,\ldots,n.
\end{equation}

The probabilities must satisfy the normalization constraint

\begin{equation}
    \sum_{i=1}^{n}p_i=1.
\end{equation}

As an example, suppose that the mean state is also constrained to be

\begin{equation}
    \sum_{i=1}^{n}ip_i=4.5.
\end{equation}

We want to find the distribution $\{p_i\}$ that maximizes

\begin{equation}
    S=-\sum_{i=1}^{n}p_i\ln p_i
\end{equation}

subject to these constraints.

Introduce Lagrange multipliers $\alpha$ and $\beta$ and define

\begin{equation}
    \mathcal{L}=-\sum_{i=1}^{n}p_i\ln p_i-\alpha\left(\sum_{i=1}^{n}p_i-1\right)-\beta\left(\sum_{i=1}^{n}ip_i-4.5\right).
\end{equation}

The stationary condition for every $p_i$ is

\begin{equation}
    \frac{\partial\mathcal{L}}{\partial p_i}=0.
\end{equation}

This gives

\begin{equation}
    -\ln p_i-1-\alpha-\beta i=0.
\end{equation}

Therefore,

\begin{equation}
    p_i=e^{-1-\alpha-\beta i}.
\end{equation}

Equivalently,

\begin{equation}
    p_i\propto e^{-\beta i}.
\end{equation}

The remaining constants are determined by the normalization and mean-value constraints.

\paragraph{Exercise}

Determine the Lagrange multipliers and hence the maximum-entropy distribution for a specified set of allowed states subject to

\begin{equation}
    \sum_i p_i=1
\end{equation}

and

\begin{equation}
    \sum_i ip_i=4.5.
\end{equation}


% ============================================================
\subsection{Outlook: physical approaches to biology}
% ============================================================

Physical approaches to biology are not defined only by the use of equations or quantitative measurements.
They also provide particular ways of deciding what constitutes an explanation or an understanding of a biological system.

Themes emphasized in the course include:

\begin{itemize}
    \item providing alternative definitions of what it means to understand a phenomenon;
    \item exposing hidden sameness between apparently unrelated biological systems;
    \item providing a different perspective on what is meant by a mechanism;
    \item allowing null hypotheses to be formulated explicitly and quantitatively;
    \item forcing us to be clear about what is known and what is not known, in the spirit of ``militant ignorance'';
    \item emphasizing conceptual knowledge rather than the accumulation of isolated facts.
\end{itemize}

A recurring theme throughout the course is that the choice of representation, model, or null hypothesis is itself part of the scientific reasoning.
The aim is not merely to calculate a quantity, but to identify a description in which the essential structure of the biological problem becomes visible.